% !PRINTED 107
\weberpage{107}
Andererseits ergibt sich aus den Formeln (8) des § 21, wenn dort $v = -\dfrac{\omega}{4}$ und $= -\dfrac{1}{4}$ gesetzt wird,
% [sic] Die Quelle läßt vor dem zweiten Gleichheitszeichen die Variable aus; der Zusammenhang legt nochmals v nahe.
\[
\begin{aligned}
\vartheta_{11}\left(\frac{\omega}{4}\right) &= i\vartheta_{01}\left(\frac{\omega}{4}\right) \\
\vartheta_{11}\left(\frac{1}{4}\right) &= \vartheta_{10}\left(\frac{1}{4}\right)
\end{aligned}
\]
und demnach sind die beiden Funktionen (12), die linear durch zwei $\vartheta$-Quadrate ausdrückbar sind, von konstanten Faktoren abgesehen, übereinstimmend mit
\[
\vartheta_{01}^2(u) + \vartheta_{11}^2(u), \quad \vartheta_{10}^2(u) - \vartheta_{11}^2(u).
\]
Die konstanten Faktoren ergeben sich unmittelbar durch $u = 0$ aus den Relationen (6), (7):
\weq{13}{\vartheta_{00}\left(0, \frac{\omega}{2}\right)\vartheta_{01}\left(u, \frac{\omega}{2}\right) = \vartheta_{01}^2(u) + \vartheta_{11}^2(u),}
\weq{14}{2\vartheta_{00}(0, 2\omega)\vartheta_{10}(2u, 2\omega) = \vartheta_{10}^2(u) - \vartheta_{11}^2(u).}
Daraus erhält man die beiden letzten Formeln, wenn man $u$ in $u + \dfrac{1}{2}$ und $u + \dfrac{\omega}{2}$ verwandelt [oder auch auf demselben Wege wie (13), (14)]:
\weq{15}{\vartheta_{00}\left(0, \frac{\omega}{2}\right)\vartheta_{00}\left(u, \frac{\omega}{2}\right) = \vartheta_{00}^2(u) + \vartheta_{10}^2(u),}
\weq{16}{2\vartheta_{00}(0, 2\omega)\vartheta_{00}(2u, 2\omega) = \vartheta_{00}^2(u) + \vartheta_{01}^2(u).}
Hieraus lassen sich mannigfache Relationen zwischen den Nullwerten der $\vartheta$-Funktionen herleiten, von denen wir nur die drei folgenden anführen, deren beide ersten aus (8), (9) fließen, während sich die letzte aus der ersten Gleichung (11) ergibt, wenn man $\omega$ durch $\omega{:}2$ ersetzt und $u = 1/4$ annimmt und berücksichtigt, daß $\vartheta_{11}(1/4) = \vartheta_{10}(1/4)$ ist.
\weq{17}{\begin{aligned}
\sqrt{\vartheta_{00}\vartheta_{01}} &= \vartheta_{01}(0, 2\omega), \\
\sqrt{\vartheta_{00}\vartheta_{10}} &= \frac{1}{\sqrt{2}}\vartheta_{10}\left(0, \frac{\omega}{2}\right), \\
\sqrt{\vartheta_{01}\vartheta_{10}} &= \vartheta_{10}\left(\frac{1}{4}, \frac{\omega}{2}\right).
\end{aligned}}
Diese Formeln sind darum von Interesse, weil sie die Quadratwurzeln als eindeutige Funktionen von $\omega$ darstellen.
